\documentclass[12pt,a4paper]{article}
\usepackage{../latex/dsdlab-report}

\begin{document}
	\reporttitlepage{۷}{طراحی \lr{UART}}

	\section*{۷.۱ مقدمه و اهداف}
	هدف این آزمایش، طراحی یک فرستنده‌گیرندهٔ ناهم‌گام همه‌منظوره\LTRfootnote{Universal Asynchronous Receiver Transmitter (UART)} است. مطابق دستورکار، بخش فرستنده هر بار یک کد هفت‌بیتی \lr{ASCII} را به‌صورت سریال می‌فرستد: نخست یک بیت شروع\LTRfootnote{Start Bit}، سپس یک بیت توازن\LTRfootnote{Parity Bit} و پس از آن هفت بیت داده؛ در پایان نیز دست‌کم یک بیت خاتمه\LTRfootnote{Stop Bit} ارسال می‌شود که در مجموع ده بیت می‌شود. بخش گیرنده پس از دریافت بیت شروع، هشت بیت مربوط به داده و توازن را به‌صورت سریال دریافت و در یک ثبات هشت‌بیتی ذخیره می‌کند. سرانجام باید با اتّصال دو واحد به یکدیگر نشان داده شود که مبادلهٔ داده به‌درستی انجام می‌شود.

	نکتهٔ مهمی که پیش از هر چیز باید روشن شود، ترتیب غیرمتعارف بیت‌های قاب\LTRfootnote{Frame} در این دستورکار است. در بیشتر پیاده‌سازی‌های متداول \lr{UART}، بیت توازن \emph{پس از} بیت‌های داده و بی‌درنگ پیش از بیت خاتمه فرستاده می‌شود؛ امّا دستورکار این آزمایش صراحتاً بیت توازن را \emph{پیش از} بیت‌های داده قرار داده است. طرح حاضر دقیقاً از همین ترتیب خواسته‌شده پیروی می‌کند و این انتخاب در سراسر فرستنده، گیرنده و میز آزمون رعایت شده است.

	\subsection*{روند انجام آزمایش}
	ابتدا ساختار قاب و ترتیب بیت‌ها از دستورکار استخراج و به‌صورت جدول درآمد. سپس فرستنده و گیرنده هر یک به‌صورت یک ماشین حالت متناهی\LTRfootnote{Finite State Machine} جداگانه نوشته شدند و یک پیمانهٔ بالا این دو را با اتّصال حلقهٔ بازگشتی\LTRfootnote{Loopback} به هم بست. برای ارزیابی، یک میز آزمون\LTRfootnote{Testbench} خودآزما\LTRfootnote{Self-Checking Testbench} ساخته شد که افزون بر مبادلهٔ درست داده، توانایی گیرنده در تشخیص خطای توازن و خطای بیت خاتمه و نیز رفتار فرستنده در برابر درخواست هم‌زمان با اشغال بودن را می‌آزماید. پس از تأیید درستی در \lr{ModelSim}، همان سه پروندهٔ سنتز‌پذیر در \lr{Quartus II} همگردانی شدند و نمای سطح انتقال ثبات‌ها\LTRfootnote{Register Transfer Level (RTL) View}، نماهای فناوری\LTRfootnote{Technology Map Viewer}، گزارش مصرف منابع و تحلیل زمانی طرح بررسی شد.

	\section*{۷.۲ ساختار قاب و تصمیم‌های طراحی}
	\subsection*{الف) ترتیب بیت‌ها}
	جدول زیر ده بیت قاب را به همان ترتیبی که روی خط ظاهر می‌شوند نشان می‌دهد.

	\begin{table}[H]
		\centering
		\small
		\begin{tabular}{|c|l|l|}
			\hline
			\textbf{شماره} & \textbf{بیت} & \textbf{مقدار} \\ \hline
			\(1\)  & شروع        & همواره صفر \\ \hline
			\(2\)  & توازن       & توازن زوج هفت بیت داده \\ \hline
			\(3\)  & دادهٔ \(D_0\) & کم‌ارزش‌ترین بیت \\ \hline
			\(4\)  & دادهٔ \(D_1\) & --- \\ \hline
			\(5\)  & دادهٔ \(D_2\) & --- \\ \hline
			\(6\)  & دادهٔ \(D_3\) & --- \\ \hline
			\(7\)  & دادهٔ \(D_4\) & --- \\ \hline
			\(8\)  & دادهٔ \(D_5\) & --- \\ \hline
			\(9\)  & دادهٔ \(D_6\) & پرارزش‌ترین بیت \\ \hline
			\(10\) & خاتمه       & همواره یک \\ \hline
		\end{tabular}
		\caption{ترتیب ده بیت قاب مطابق دستورکار}
	\end{table}

	بیت‌های داده از کم‌ارزش به پرارزش فرستاده می‌شوند، که قرارداد متعارف \lr{UART} است. بیت توازن نیز از نوع زوج\LTRfootnote{Even Parity} است و در کد با عملگر کاهشی یای‌انحصاری روی هفت بیت داده ساخته می‌شود؛ یعنی اگر شمار یک‌های داده زوج باشد بیت توازن صفر و در غیر این صورت یک می‌شود.

	\subsection*{ب) خط بی‌کار و آشکارسازی شروع}
	خط ارسال در حالت بی‌کار روی یک نگه داشته می‌شود. چون بیت شروع همواره صفر است، گیرنده تنها با دیدن یک سطح پایین روی خط بی‌کارِ یک، آغاز قاب را تشخیص می‌دهد. این همان سازوکار متعارف \lr{UART} است و به همین دلیل نیازی به هیچ سیگنال هم‌گام‌سازی جداگانه‌ای میان فرستنده و گیرنده نیست؛ صفت «ناهم‌گام» در نام این واحد نیز به همین ویژگی اشاره دارد.

	\subsection*{ج) طول هر بیت و پارامتر \lr{TICKS\_BIT}}
	در یک \lr{UART} واقعی، طول هر بیت روی خط از نرخ باود\LTRfootnote{Baud Rate} تعیین می‌شود و معمولاً بسیار بزرگ‌تر از یک پالس ساعت است. در این طرح، طول هر بیت با پارامتر \lr{TICKS\_BIT} بر حسب شمار پالس‌های ساعت بیان شده است و مقدار پیش‌فرض آن دو است. با این کار، هم شبیه‌سازی در زمان معقول اجرا می‌شود و هم طرح به یک نرخ باود مشخّص وابسته نمی‌ماند؛ برای رسیدن به نرخ باود دلخواه کافی است این پارامتر برابر نسبت بسامد ساعت به نرخ باود گرفته شود.

	هر دو پیمانهٔ فرستنده و گیرنده این پارامتر را می‌پذیرند و پیمانهٔ بالا آن را به هر دو منتقل می‌کند، بنابراین دو سوی خط همواره با یک طول بیت یکسان کار می‌کنند.

	\subsection*{د) وارسی درستی قاب در گیرنده}
	دستورکار تنها ذخیرهٔ داده و توازن در یک ثبات هشت‌بیتی را خواسته است. در این طرح، افزون بر ذخیرهٔ داده، خروجی سومی نیز به نام \lr{data\_correct} افزوده شد که سلامت قاب دریافتی را گزارش می‌کند. این خروجی تنها هنگامی یک می‌شود که هر دو شرط زیر برقرار باشند: بیت خاتمه در جای خود برابر یک باشد و توازن محاسبه‌شده از روی هفت بیت داده با بیت توازن دریافتی بخواند. به این ترتیب، هم خطای تک‌بیتی در داده و هم از دست رفتن هم‌گامی قاب آشکار می‌شود.

	\section*{۷.۳ توصیف فرستنده}
	فرستنده یک ماشین حالت با پنج حالت است: بی‌کار، شروع، توازن، داده و خاتمه. در حالت بی‌کار، خط روی یک نگه داشته می‌شود و پرچم اشغال صفر است. با یک شدن \lr{data\_new}، هفت بیت داده در یک ثبات درونی نگه داشته می‌شود، بیت توازن همان‌جا محاسبه و ذخیره می‌گردد، خط به صفر می‌رود و ماشین وارد حالت شروع می‌شود.

	نکتهٔ درخور توجّه آن است که درخواست ارسال تنها در حالت بی‌کار پذیرفته می‌شود. اگر در میانهٔ ارسال یک قاب، درخواست تازه‌ای برسد، آن درخواست به‌سادگی نادیده گرفته می‌شود و قاب در حال ارسال آسیب نمی‌بیند. این تصمیم در بخش شبیه‌سازی به‌صورت جداگانه آزموده شده است.

	در هر یک از حالت‌های بعدی، یک شمارنده به اندازهٔ \lr{TICKS\_BIT} پالس ساعت می‌شمارد و سپس بیت بعدی را روی خط می‌گذارد. در حالت داده، شمارندهٔ \lr{bit\_index} از صفر تا شش پیش می‌رود و پس از فرستادن \(D_6\)، خط به یک بازمی‌گردد و ماشین وارد حالت خاتمه می‌شود. پرچم اشغال تا پایان بیت خاتمه بالا می‌ماند تا میز آزمون یا هر مصرف‌کنندهٔ دیگری بداند که فرستنده هنوز آزاد نیست.

	\vcode{code/UARTSender.v}{فرستنده‌ی \lr{UART}}

	\section*{۷.۴ توصیف گیرنده}
	گیرنده چهار حالت دارد: بی‌کار، شروع، داده و خاتمه. در حالت بی‌کار، تنها پایین رفتن خط دنبال می‌شود. پس از آشکارسازی شروع، ماشین به اندازهٔ یک طول بیت صبر می‌کند و سپس بیت توازن را نمونه‌برداری می‌کند؛ پس از آن هفت بیت داده یکی‌یکی در ثبات درونی نوشته می‌شوند و سرانجام در حالت خاتمه، بیت خاتمه بررسی می‌گردد.

	زمان‌بندی نمونه‌برداری در اینجا اهمّیّت دارد. چون شمارش هر بیت از لحظهٔ آشکارسازی لبهٔ شروع آغاز می‌شود و نمونه‌برداری در پایان شمارش انجام می‌گیرد، هر بیت در میانهٔ بازهٔ خود و نه در لبه‌های آن خوانده می‌شود. همین فاصله از لبه‌هاست که خواندن را در برابر اختلاف جزئی زمان‌بندی میان دو سوی خط مقاوم می‌کند.

	خروجی \lr{data\_new\_rec} تنها برای یک پالس ساعت بالا می‌رود و پایان یک قاب کامل را اعلام می‌کند؛ حال آنکه \lr{data\_rec} و \lr{data\_correct} تا رسیدن قاب بعدی مقدار خود را نگه می‌دارند. این تفکیک باعث می‌شود مصرف‌کنندهٔ خروجی بتواند با یک پالس، لحظهٔ آمادگی دادهٔ تازه را تشخیص دهد و در عین حال داده تا استفادهٔ بعدی پایدار بماند.

	\vcode{code/UARTReceiver.v}{گیرنده‌ی \lr{UART}}

	\section*{۷.۵ اتّصال دو واحد}
	دستورکار خواسته است که با اتّصال دو واحد به یکدیگر، درستی مبادلهٔ داده نشان داده شود. پیمانهٔ بالا دقیقاً همین کار را انجام می‌دهد: یک نمونه از فرستنده و یک نمونه از گیرنده می‌سازد و خروجی \lr{tx} فرستنده را مستقیماً به ورودی \lr{rx} گیرنده می‌بندد. این همان آرایش حلقهٔ بازگشتی است و در آن، هر قابی که فرستنده روی خط می‌گذارد بی‌درنگ توسط گیرنده خوانده و بازسازی می‌شود.

	\vcode{code/UARTTop.v}{پیمانه‌ی بالا و اتّصال حلقه‌ی بازگشتی فرستنده به گیرنده}

	نمای سطح انتقال ثبات‌های پیمانهٔ بالا همین ساختار را نشان می‌دهد: دو بلوک مجزّای فرستنده و گیرنده که تنها پیوند میانشان، خط سریال است.

	\shot{screenshots/rtl_top.png}{نمای سطح انتقال ثبات‌های پیمانه‌ی بالا و مسیر حلقه‌ی بازگشتی}

	\section*{۷.۶ میز آزمون}
	میز آزمون افزون بر پیمانهٔ بالا، یک گیرندهٔ مستقل دیگر نیز می‌سازد که ورودی سریال آن مستقیماً از میز آزمون تغذیه می‌شود. هدف از این گیرندهٔ دوم، تزریق قاب‌های معیوب است: چون فرستندهٔ سالم هرگز قاب نادرست تولید نمی‌کند، برای آزمودن توانایی تشخیص خطا باید قاب را دستی و بیت‌به‌بیت ساخت. وظیفهٔ \lr{send\_manual\_frame} همین کار را می‌کند و دو ورودی جداگانه دارد که با آن‌ها می‌توان بیت توازن یا بیت خاتمه را عمداً خراب کرد.

	آزمون از پنج مرحله تشکیل شده است:
	\begin{itemize}
		\item \textbf{مرحلهٔ ۱ --- قاب سالم}: یک کد هفت‌بیتی از مسیر حلقهٔ بازگشتی فرستاده و بررسی می‌شود که همان مقدار در گیرنده بازسازی شود و پرچم درستی بالا برود.
		\item \textbf{مرحلهٔ ۲ --- خطای توازن}: قابی با بیت توازن وارونه ساخته می‌شود و بررسی می‌گردد که گیرنده آن را نادرست اعلام کند.
		\item \textbf{مرحلهٔ ۳ --- خطای بیت خاتمه}: قابی با بیت خاتمهٔ صفر ساخته می‌شود و انتظار می‌رود پرچم درستی پایین بماند.
		\item \textbf{مرحلهٔ ۴ --- دو قاب پیاپی}: دو کد متفاوت پشت سر هم فرستاده می‌شوند تا نشان داده شود ماشین پس از هر قاب به‌درستی به حالت بی‌کار بازمی‌گردد.
		\item \textbf{مرحلهٔ ۵ --- درخواست هم‌زمان با اشغال}: در میانهٔ ارسال یک قاب، درخواست تازه‌ای اعمال می‌شود. آزمون هم بررسی می‌کند که قاب در حال ارسال دست‌نخورده بماند و هم اینکه از درخواست نادیده‌گرفته‌شده هیچ قاب اضافه‌ای ساخته نشود.
	\end{itemize}

	مرحلهٔ پنجم دو وارسی جداگانه دارد و همین آن را از یک آزمون سطحی متمایز می‌کند. وارسی نخست مقدار بازسازی‌شده را می‌سنجد و وارسی دوم، پس از آزاد شدن فرستنده، به مدّت ده طول بیت دیگر منتظر می‌ماند تا مطمئن شود هیچ پالس \lr{data\_new\_rec} تازه‌ای ظاهر نمی‌شود.

	\vcode{code/Tester.v}{میز آزمون خودآزمای \lr{UART}}

	\section*{۷.۷ اجرا در \lr{ModelSim}}
	اجرای شبیه‌سازی با یک پروندهٔ فرمان انجام می‌شود که کتابخانه را از نو می‌سازد، چهار پرونده را همگردانی می‌کند، سیگنال‌ها را در دو دستهٔ حلقهٔ بازگشتی و گیرندهٔ آزمون خطا به پنجرهٔ موج می‌افزاید و شبیه‌سازی را تا پایان اجرا می‌کند.

	\docode{code/run.do}{پرونده‌ی فرمان اجرای شبیه‌سازی}

	خروجی متنی آزمون نشان می‌دهد که هر پنج مرحله با موفّقیّت گذشته‌اند. تصویرهای زیر رونوشت کامل اجرا را از آغاز تا پیام پایانی دنبال می‌کنند.

	\shotw{screenshots/transcript_1.png}{0.82}{رونوشت اجرا؛ مرحله‌های یک تا سه و آغاز مرحله‌ی چهارم}
	\shotw{screenshots/transcript_2.png}{0.82}{ادامه‌ی رونوشت؛ مرحله‌های چهارم و پنجم}
	\shotw{screenshots/transcript_3.png}{0.82}{پیام پایانی و گزارش قبولی همه‌ی آزمون‌ها}

	آخرین سطرها نشان می‌دهند که شمارندهٔ خطا صفر مانده و شبیه‌سازی در لحظهٔ ۱۶۹۶ نانوثانیه به پایان رسیده است.

	\section*{۷.۸ تحلیل شکل موج‌ها}
	\subsection*{الف) کلید خواندن شکل موج}
	پیش از هر تحلیلی باید دو عدد را در دست داشت، چون همهٔ بازه‌های زمانی این آزمایش از آن‌ها به دست می‌آیند. نخست، دورهٔ ساعت میز آزمون ۱۰ نانوثانیه است. دوم، پارامتر \lr{TICKS\_BIT} برابر دو است، یعنی هر بیت روی خط به اندازهٔ دو پالس ساعت پایدار می‌ماند. از ضرب این دو نتیجه می‌شود:
	\[
	\text{طول هر بیت} = 2 \times 10 = 20\ \text{نانوثانیه},
	\qquad
	\text{طول یک قاب کامل} = 10 \times 20 = 200\ \text{نانوثانیه}.
	\]

	همین دو عدد، نکتهٔ کلیدی ظاهر شکل موج را نیز توضیح می‌دهند. در نگاه نخست ممکن است پرسیده شود چرا پالس‌های روی خط \lr{tx} پهنای یکسانی ندارند و برخی دو برابر برخی دیگرند. پاسخ این است که خط \lr{tx} هیچ نشانه‌ای برای مرز میان بیت‌ها ندارد؛ اگر دو بیت پیاپی مقدار یکسانی داشته باشند، روی خط به یک پالس پیوستهٔ ۴۰ نانوثانیه‌ای دیده می‌شوند. بنابراین پهنای هر پالس همواره مضرب صحیحی از ۲۰ نانوثانیه است و شمار آن مضرب، تعداد بیت‌های هم‌مقدار پشت سر هم را می‌گوید. همین ویژگی است که ایجاب می‌کند گیرنده به‌جای شمردن لبه‌ها، از لحظهٔ آشکارسازی شروع، زمان را خودش بشمارد.

	\subsection*{ب) سیگنال‌های پنجره‌ی موج}
	جدول زیر سیگنال‌های نمایش‌داده‌شده در پنجرهٔ موج و معنای آن‌ها را گردآوری می‌کند.

	\begin{table}[H]
		\centering
		\small
		\begin{tabular}{|l|l|l|}
			\hline
			\textbf{دسته} & \textbf{سیگنال} & \textbf{معنا} \\ \hline
			حلقهٔ بازگشتی & \lr{clk}          & ساعت با دورهٔ ۱۰ نانوثانیه \\ \hline
			حلقهٔ بازگشتی & \lr{rstN}         & بازنشانی ناهم‌زمان فعّال‌صفر \\ \hline
			حلقهٔ بازگشتی & \lr{new\_data}    & درخواست ارسال یک کد تازه \\ \hline
			حلقهٔ بازگشتی & \lr{send\_data}   & کد هفت‌بیتی ورودی فرستنده \\ \hline
			حلقهٔ بازگشتی & \lr{tx}           & خط سریال خروجی فرستنده \\ \hline
			حلقهٔ بازگشتی & \lr{busy}         & اشغال بودن فرستنده \\ \hline
			حلقهٔ بازگشتی & \lr{rec\_data}    & کد بازسازی‌شده در گیرنده \\ \hline
			حلقهٔ بازگشتی & \lr{rec\_new\_data} & پالس یک‌کلاکی پایان قاب \\ \hline
			حلقهٔ بازگشتی & \lr{correct\_data} & سلامت قاب دریافتی \\ \hline
			آزمون خطا & \lr{fault\_rx}        & خط سریالی که میز آزمون دستی می‌راند \\ \hline
			آزمون خطا & \lr{fault\_data}      & دادهٔ گیرنده‌ی آزمون خطا \\ \hline
			آزمون خطا & \lr{fault\_new\_data} & پالس پایان قاب معیوب \\ \hline
			آزمون خطا & \lr{fault\_correct}   & نتیجه‌ی وارسی قاب معیوب \\ \hline
		\end{tabular}
		\caption{سیگنال‌های پنجره‌ی موج و معنای آن‌ها}
	\end{table}

	\shot{screenshots/wave_full.png}{شکل موج کامل شبیه‌سازی و هر پنج مرحله‌ی آزمون}

	در نمای کامل، پنج بازهٔ فعّالیّت از هم جدا دیده می‌شوند که به‌ترتیب به پنج مرحلهٔ آزمون مربوط‌اند. جدول زیر این بازه‌ها را مشخّص می‌کند تا در ادامه بتوان به هر یک ارجاع داد.

	\begin{table}[H]
		\centering
		\small
		\begin{tabular}{|c|c|l|}
			\hline
			\textbf{مرحله} & \textbf{بازه (نانوثانیه)} & \textbf{آنچه در آن رخ می‌دهد} \\ \hline
			\(1\) & \(65\)--\(265\)   & قاب سالم روی مسیر حلقهٔ بازگشتی \\ \hline
			\(2\) & \(280\)--\(480\)  & قاب دستی با بیت توازن وارونه \\ \hline
			\(3\) & \(520\)--\(720\)  & قاب دستی با بیت خاتمهٔ صفر \\ \hline
			\(4\) & \(805\)--\(1215\) & دو قاب سالم پیاپی \\ \hline
			\(5\) & \(1225\)--\(1425\) & قاب سالم به‌همراه درخواست نادیده‌گرفته‌شده \\ \hline
		\end{tabular}
		\caption{بازه‌های زمانی پنج مرحله‌ی آزمون در شکل موج کامل}
	\end{table}

	\subsection*{ج) تحلیل بازه‌به‌بازه‌ی قاب سالم}
	مرحلهٔ نخست، کد \(1011001_2\) را می‌فرستد. برای فهمیدن اینکه چرا خط \lr{tx} دقیقاً همان شکلی می‌شود که در نما دیده می‌شود، ابتدا باید بیت‌های این کد را جدا کرد. بیت‌ها از کم‌ارزش به پرارزش عبارت‌اند از
	\[
	D_0=1,\; D_1=0,\; D_2=0,\; D_3=1,\; D_4=1,\; D_5=0,\; D_6=1,
	\]
	یعنی چهار بیت یک و سه بیت صفر. چون شمار یک‌ها زوج است، بیت توازن زوج برابر صفر می‌شود:
	\[
	P = D_0\oplus D_1\oplus\cdots\oplus D_6 = 1\oplus0\oplus0\oplus1\oplus1\oplus0\oplus1 = 0 .
	\]

	درخواست ارسال در لحظهٔ ۶۰ نانوثانیه اعمال می‌شود و فرستنده آن را در نخستین لبهٔ فعّال پس از آن، یعنی ۶۵ نانوثانیه، می‌پذیرد. از همین‌جا قاب آغاز می‌شود. جدول زیر هر ده بیت قاب را با بازهٔ زمانی دقیق و دلیل مقدار آن می‌آورد.

	\begin{table}[H]
		\centering
		\footnotesize
		\begin{tabular}{|c|c|c|l|}
			\hline
			\textbf{بازه (نانوثانیه)} & \textbf{بیت} & \textbf{\lr{tx}} & \textbf{چرا این مقدار} \\ \hline
			\(65\)--\(85\)   & شروع        & \(0\) & قرارداد قاب؛ پایین رفتن خط بی‌کار، آغاز را اعلام می‌کند \\ \hline
			\(85\)--\(105\)  & توازن       & \(0\) & شمار یک‌های داده زوج است \\ \hline
			\(105\)--\(125\) & \(D_0\)     & \(1\) & کم‌ارزش‌ترین بیت کد \\ \hline
			\(125\)--\(145\) & \(D_1\)     & \(0\) & --- \\ \hline
			\(145\)--\(165\) & \(D_2\)     & \(0\) & --- \\ \hline
			\(165\)--\(185\) & \(D_3\)     & \(1\) & --- \\ \hline
			\(185\)--\(205\) & \(D_4\)     & \(1\) & --- \\ \hline
			\(205\)--\(225\) & \(D_5\)     & \(0\) & --- \\ \hline
			\(225\)--\(245\) & \(D_6\)     & \(1\) & پرارزش‌ترین بیت کد \\ \hline
			\(245\)--\(265\) & خاتمه       & \(1\) & قرارداد قاب؛ بازگشت خط به حالت بی‌کار \\ \hline
		\end{tabular}
		\caption{تحلیل بازه‌به‌بازه‌ی قاب سالم مرحله‌ی نخست}
	\end{table}

	اکنون می‌توان به پرسش اصلی پاسخ داد: چرا شکل موج \lr{tx} در این بازه همان است که دیده می‌شود؟ اگر بیت‌های هم‌مقدار پشت سر هم را با هم ادغام کنیم، خط \lr{tx} تنها در شش لحظه تغییر حالت می‌دهد و نه ده بار. جدول زیر همین را نشان می‌دهد و دقیقاً با لحظه‌های تغییری که در شکل موج دیده می‌شوند می‌خواند.

	\begin{table}[H]
		\centering
		\footnotesize
		\begin{tabular}{|c|c|c|l|}
			\hline
			\textbf{بازه (نانوثانیه)} & \textbf{مقدار} & \textbf{پهنا} & \textbf{ترکیب بیت‌ها} \\ \hline
			\(65\)--\(105\)  & \(0\) & \(40\) & شروع و توازن، هر دو صفر \\ \hline
			\(105\)--\(125\) & \(1\) & \(20\) & تنها \(D_0\) \\ \hline
			\(125\)--\(165\) & \(0\) & \(40\) & \(D_1\) و \(D_2\)، هر دو صفر \\ \hline
			\(165\)--\(205\) & \(1\) & \(40\) & \(D_3\) و \(D_4\)، هر دو یک \\ \hline
			\(205\)--\(225\) & \(0\) & \(20\) & تنها \(D_5\) \\ \hline
			\(225\)--\(265\) & \(1\) & \(40\) & \(D_6\) و بیت خاتمه، هر دو یک \\ \hline
		\end{tabular}
		\caption{پالس‌های دیده‌شده روی خط \lr{tx} و ترکیب بیتی هر یک}
	\end{table}

	پس آنچه در نما به‌صورت یک پالس پهن دیده می‌شود، دو بیت جداگانه با مقدار یکسان است و نه یک بیت طولانی. توجّه به همین نکته، تفاوت میان خواندن درست و نادرست شکل موج \lr{UART} است.

	\shot{screenshots/wave_zoom_1.png}{نمای بزرگ‌شده‌ی مرحله‌ی نخست؛ قاب سالم روی خط \lr{tx} و بازسازی آن در گیرنده}

	\subsection*{د) لحظه‌های نمونه‌برداری گیرنده}
	تا اینجا رفتار فرستنده توضیح داده شد. برای اینکه روشن شود چرا خروجی گیرنده دقیقاً همان مقدار ارسالی می‌شود، باید دید گیرنده در چه لحظه‌هایی خط را می‌خواند. گیرنده در لحظهٔ ۷۵ نانوثانیه، یعنی نخستین لبهٔ فعّال پس از پایین رفتن خط، شروع را آشکار می‌کند و از همان‌جا شمارش خود را آغاز می‌نماید. جدول زیر لحظهٔ خواندن هر بیت را می‌آورد.

	\begin{table}[H]
		\centering
		\footnotesize
		\begin{tabular}{|c|c|c|c|}
			\hline
			\textbf{لحظه (نانوثانیه)} & \textbf{چه خوانده می‌شود} & \textbf{بازه‌ی آن بیت} & \textbf{مقدار} \\ \hline
			\(75\)  & آشکارسازی شروع & \(65\)--\(85\)   & \(0\) \\ \hline
			\(95\)  & بیت توازن      & \(85\)--\(105\)  & \(0\) \\ \hline
			\(115\) & \(D_0\)        & \(105\)--\(125\) & \(1\) \\ \hline
			\(135\) & \(D_1\)        & \(125\)--\(145\) & \(0\) \\ \hline
			\(155\) & \(D_2\)        & \(145\)--\(165\) & \(0\) \\ \hline
			\(175\) & \(D_3\)        & \(165\)--\(185\) & \(1\) \\ \hline
			\(195\) & \(D_4\)        & \(185\)--\(205\) & \(1\) \\ \hline
			\(215\) & \(D_5\)        & \(205\)--\(225\) & \(0\) \\ \hline
			\(235\) & \(D_6\)        & \(225\)--\(245\) & \(1\) \\ \hline
			\(255\) & بیت خاتمه      & \(245\)--\(265\) & \(1\) \\ \hline
		\end{tabular}
		\caption{لحظه‌های نمونه‌برداری گیرنده و بازه‌ی متناظر هر بیت}
	\end{table}

	ستون سوم این جدول، مهم‌ترین نکتهٔ زمان‌بندی طرح را آشکار می‌کند: هر لحظهٔ نمونه‌برداری دقیقاً ۱۰ نانوثانیه پس از آغاز بازهٔ بیت خود قرار دارد، یعنی درست در میانهٔ آن بازهٔ ۲۰ نانوثانیه‌ای. بنابراین در هر سو ۱۰ نانوثانیه فاصله تا لبه‌های بیت باقی می‌ماند و همین فاصله است که خواندن را در برابر اختلاف جزئی ساعت دو سوی خط ایمن می‌کند. اگر نمونه‌برداری نزدیک لبه انجام می‌شد، کوچک‌ترین ناهم‌خوانی می‌توانست بیت همسایه را به‌جای بیت درست بخواند.

	\subsection*{ه) چرا خروجی گیرنده در لحظه‌ی ۲۵۵ نانوثانیه ظاهر می‌شود}
	در شکل موج دیده می‌شود که \lr{rec\_data} تا لحظهٔ ۲۵۵ نانوثانیه مقدار پیشین خود را نگه می‌دارد و تنها در آن لحظه به \(1011001\) تغییر می‌کند؛ هم‌زمان \lr{rec\_new\_data} برای یک پالس ساعت بالا می‌رود و \lr{correct\_data} یک می‌شود. دلیل این زمان‌بندی از جدول بالا روشن است: گیرنده تا پیش از خواندن بیت خاتمه هیچ داوری‌ای نمی‌کند، چون تا آن لحظه هنوز نمی‌داند قاب سالم بوده است یا نه. در لحظهٔ ۲۵۵ نانوثانیه سه کار هم‌زمان انجام می‌شود:
	\begin{itemize}
		\item محتوای ثبات درونی به خروجی \lr{rec\_data} منتقل می‌شود.
		\item پرچم \lr{rec\_new\_data} برای یک پالس بالا می‌رود تا پایان یک قاب کامل اعلام شود.
		\item دو شرط سلامت با هم سنجیده می‌شوند: بیت خاتمهٔ خوانده‌شده باید یک باشد و توازن محاسبه‌شده از روی هفت بیت ذخیره‌شده باید با بیت توازن دریافتی بخواند. در این قاب هر دو برقرارند، پس \lr{correct\_data} یک می‌شود.
	\end{itemize}

	از اینجا تأخیر کلّ مسیر نیز به دست می‌آید: از لحظهٔ ۶۵ که فرستنده درخواست را پذیرفت تا لحظهٔ ۲۵۵ که داده در خروجی گیرنده حاضر شد، ۱۹۰ نانوثانیه گذشته است. این عدد ده نانوثانیه کمتر از طول کامل قاب است، چون خروجی در لحظهٔ \emph{نمونه‌برداری} بیت خاتمه به‌روز می‌شود و نه در پایان آن بیت.

	پرچم \lr{busy} نیز رفتار متناظری دارد: از ۶۵ تا ۲۶۵ نانوثانیه بالا می‌ماند، یعنی دقیقاً ۲۰۰ نانوثانیه که برابر طول کامل قاب است. تفاوت ۱۰ نانوثانیه‌ای میان لحظهٔ پایین آمدن \lr{busy} و لحظهٔ ظاهر شدن خروجی گیرنده، همان اختلافی است که در بند پیش توضیح داده شد.

	\subsection*{و) قاب معیوب توازن}
	در مرحلهٔ دوم، میز آزمون همان کد \(1011001_2\) را می‌فرستد امّا بیت توازن را عمداً وارونه می‌کند؛ یعنی به‌جای مقدار درست صفر، مقدار یک روی خط می‌گذارد. جدول زیر بازه‌های این قاب را نشان می‌دهد.

	\begin{table}[H]
		\centering
		\footnotesize
		\begin{tabular}{|c|c|c|l|}
			\hline
			\textbf{بازه (نانوثانیه)} & \textbf{بیت} & \textbf{\lr{fault\_rx}} & \textbf{وضعیّت} \\ \hline
			\(280\)--\(300\) & شروع  & \(0\) & درست \\ \hline
			\(300\)--\(320\) & توازن & \(1\) & \textbf{وارونه‌ی مقدار درست} \\ \hline
			\(320\)--\(460\) & \(D_0\) تا \(D_6\) & \(1001101\) & درست \\ \hline
			\(460\)--\(480\) & خاتمه & \(1\) & درست \\ \hline
		\end{tabular}
		\caption{بازه‌های قاب معیوب توازن در مرحله‌ی دوم}
	\end{table}

	گیرنده این قاب را بی‌هیچ مشکلی تا انتها دریافت می‌کند و در لحظهٔ ۴۶۵ نانوثانیه پالس \lr{fault\_new\_data} را می‌دهد؛ یعنی از دید ساختار قاب، همه‌چیز عادی بوده است. امّا در همان لحظه، توازن محاسبه‌شده از روی هفت بیت داده صفر به دست می‌آید و با بیت توازن دریافتی که یک است نمی‌خواند، پس \lr{fault\_correct} صفر می‌ماند. این دقیقاً همان رفتاری است که انتظار می‌رفت: قاب کامل شد، ولی سالم نبود.

	\subsection*{ز) قاب معیوب بیت خاتمه}
	در مرحلهٔ سوم، بیت توازن درست است ولی بیت خاتمه به‌جای یک، صفر گذاشته شده است.

	\begin{table}[H]
		\centering
		\footnotesize
		\begin{tabular}{|c|c|c|l|}
			\hline
			\textbf{بازه (نانوثانیه)} & \textbf{بیت} & \textbf{\lr{fault\_rx}} & \textbf{وضعیّت} \\ \hline
			\(520\)--\(540\) & شروع  & \(0\) & درست \\ \hline
			\(540\)--\(560\) & توازن & \(0\) & درست \\ \hline
			\(560\)--\(700\) & \(D_0\) تا \(D_6\) & \(1001101\) & درست \\ \hline
			\(700\)--\(720\) & خاتمه & \(0\) & \textbf{باید یک می‌بود} \\ \hline
		\end{tabular}
		\caption{بازه‌های قاب معیوب بیت خاتمه در مرحله‌ی سوم}
	\end{table}

	گیرنده در لحظهٔ ۷۰۵ نانوثانیه بیت خاتمه را می‌خواند و مقدار صفر می‌بیند. چون یکی از دو شرط سلامت نقض شده است، \lr{fault\_correct} دوباره صفر می‌ماند، هرچند این بار توازن کاملاً درست بوده است. مقایسهٔ این دو مرحله نشان می‌دهد که هر یک از دو شرط به‌تنهایی می‌تواند قاب را نامعتبر کند و پرچم سلامت، یای منطقی آن دو نیست بلکه واو منطقی آن‌هاست.

	\shot{screenshots/wave_zoom_2.png}{نمای بزرگ‌شده‌ی قاب‌های معیوب و پایین ماندن پرچم درستی}

	\subsection*{ح) بازه‌ی اشغال بودن فرستنده}
	مرحلهٔ پنجم، رفتار فرستنده را در برابر درخواستی می‌سنجد که در میانهٔ ارسال می‌رسد. دنبالهٔ رویدادها چنین است:

	\begin{table}[H]
		\centering
		\footnotesize
		\begin{tabular}{|c|l|}
			\hline
			\textbf{لحظه (نانوثانیه)} & \textbf{رویداد} \\ \hline
			\(1220\) & درخواست نخست با کد \(1100101\) اعمال می‌شود \\ \hline
			\(1225\) & فرستنده درخواست را می‌پذیرد و \lr{busy} بالا می‌رود \\ \hline
			\(1260\) & درخواست دوم با کد \(0011110\) در میانه‌ی ارسال اعمال می‌شود \\ \hline
			\(1415\) & قاب نخست کامل می‌شود؛ \lr{rec\_data} برابر \(1100101\) \\ \hline
			\(1425\) & \lr{busy} پایین می‌رود و فرستنده آزاد می‌شود \\ \hline
		\end{tabular}
		\caption{دنباله‌ی رویدادهای مرحله‌ی پنجم}
	\end{table}

	در لحظهٔ ۱۲۶۰ نانوثانیه، ماشین فرستنده در میانهٔ حالت داده است و نه در حالت بی‌کار. چون در کد، پذیرش درخواست تنها در حالت بی‌کار بررسی می‌شود، این درخواست هیچ اثری نمی‌گذارد: نه ثبات دادهٔ درونی بازنویسی می‌شود و نه ماشین حالت از مسیر خود منحرف می‌گردد. شاهد این ادّعا در شکل موج آن است که با وجود تغییر \lr{send\_data} به \(0011110\) در لحظهٔ ۱۲۶۰، مقدار بازسازی‌شده در گیرنده در لحظهٔ ۱۴۱۵ همچنان \(1100101\) است؛ یعنی همان کدی که پیش از اشغال شدن فرستنده قفل شده بود.

	میز آزمون یک گام فراتر می‌رود و پس از آزاد شدن فرستنده، به مدّت ده طول بیت دیگر خط را می‌پاید تا مطمئن شود درخواست نادیده‌گرفته‌شده بعداً هم قاب اضافه‌ای نمی‌سازد. در این بازه هیچ پالس تازه‌ای روی \lr{rec\_new\_data} ظاهر نمی‌شود، پس درخواست به‌کلّی دور ریخته شده و صرفاً به تعویق نیفتاده است.

	\subsection*{ط) جمع‌بندی تحلیل}
	آنچه از بررسی این بازه‌ها به دست می‌آید چنین است:
	\begin{itemize}
		\item هر بیت روی خط دقیقاً ۲۰ نانوثانیه و هر قاب دقیقاً ۲۰۰ نانوثانیه طول می‌کشد و همهٔ لحظه‌های تغییر، مضرب صحیحی از این عددند.
		\item پهنای نایکسان پالس‌های \lr{tx} نه نشانهٔ بی‌نظمی، بلکه پیامد ادغام دیداری بیت‌های هم‌مقدار پیاپی است.
		\item گیرنده هر بیت را در میانهٔ بازهٔ آن و با ۱۰ نانوثانیه حاشیه از هر دو لبه می‌خواند.
		\item خروجی گیرنده تنها در لحظهٔ نمونه‌برداری بیت خاتمه به‌روز می‌شود، که تأخیر کلّ مسیر را ۱۹۰ نانوثانیه می‌کند.
		\item پرچم پایان قاب و پرچم سلامت قاب مستقل از یکدیگرند؛ در هر دو قاب معیوب، اوّلی بالا رفت و دوّمی پایین ماند.
		\item درخواست ارسال در بازهٔ اشغال بودن، نه قاب جاری را تغییر داد و نه بعداً قاب تازه‌ای ساخت.
	\end{itemize}

	\section*{۷.۹ سنتز و پیاده‌سازی نرم‌افزاری}
	\subsection*{الف) مصرف منابع}
	سه پروندهٔ سنتز‌پذیر طرح، بدون میز آزمون، در \lr{Quartus II 13.0.1} و برای خانوادهٔ \lr{Cyclone IV GX} همگردانی شدند. انتخاب قطعه روی حالت خودکار گذاشته شد و ابزار قطعهٔ \lr{EP4CGX15BF14C6} را برگزید. نتیجهٔ همگردانی ۲۱۵ عنصر منطقی شامل ۲۰۹ تابع ترکیبی و ۱۰۶ ثبات اختصاصی، و ۲۱ پایه است.

	\subsection*{ب) بررسی شمار ثبات‌ها}
	شمار ۱۰۶ ثبات برای مداری با این اندازه در نگاه نخست زیاد به نظر می‌رسد و ارزش بررسی دارد. گزارش سنتز، ۵۰ ثبات را به فرستنده و ۵۶ ثبات را به گیرنده نسبت می‌دهد. اگر ثبات‌های ساختاری هر پیمانه شمرده شوند، به این اعداد می‌رسیم:
	\begin{itemize}
		\item فرستنده: پنج بیت برای ثبات حالت با رمزگذاری یک‌داغ\LTRfootnote{One-Hot Encoding}، سه بیت \lr{bit\_index}، هفت بیت دادهٔ ذخیره‌شده، یک بیت توازن، و دو بیت برای خروجی‌های \lr{tx} و \lr{busy}؛ در مجموع هجده بیت.
		\item گیرنده: چهار بیت حالت، سه بیت \lr{bit\_index}، هفت بیت ثبات داخلی، یک بیت توازن دریافتی، هفت بیت خروجی داده و دو بیت برای دو پرچم؛ در مجموع بیست و چهار بیت.
	\end{itemize}
	اختلاف این اعداد با گزارش، در هر دو پیمانه دقیقاً سی و دو بیت است. منشأ آن روشن است: شمارندهٔ \lr{tick\_count} در هر دو پیمانه از نوع \lr{integer} تعریف شده و این نوع در \lr{Verilog} همواره سی و دو بیت است. بنابراین دو شمارنده که برای مقدار پیش‌فرض \lr{TICKS\_BIT} تنها به یک بیت نیاز دارند، در عمل شصت و چهار فلیپ‌فلاپ گرفته‌اند؛ یعنی نزدیک به شصت درصد کلّ ثبات‌های طرح.

	این نکته بر درستی عملکرد اثری ندارد و در شبیه‌سازی نیز دیده نمی‌شود، امّا از نظر مصرف منابع بهبودپذیر است: اگر این شمارنده به‌جای \lr{integer} با پهنای \(\lceil\log_2(\text{\lr{TICKS\_BIT}})\rceil\) بیت تعریف می‌شد، شمار ثبات‌های طرح از ۱۰۶ به حدود ۴۲ کاهش می‌یافت. آوردن این مشاهده در گزارش از آن روست که تفاوت میان «کد درست» و «کد بهینه» را نشان می‌دهد.

	\subsection*{ج) نماهای سطح انتقال ثبات‌ها و فناوری}
	نمای سطح انتقال ثبات‌های فرستنده، ثبات حالت، شمارندهٔ بیت، ثبات دادهٔ ذخیره‌شده و منطق تولید بیت توازن را نشان می‌دهد. در نمای گیرنده نیز ثبات حالت، ثبات داده و منطق مقایسهٔ توازن دیده می‌شود.

	\shot{screenshots/rtl_sender.png}{نمای سطح انتقال ثبات‌های فرستنده}
	\shot{screenshots/rtl_receiver.png}{نمای سطح انتقال ثبات‌های گیرنده}

	نماهای فناوری پس از نگاشت و پس از جانمایی نشان می‌دهند که این ساختار به جدول‌های جست‌وجو و فلیپ‌فلاپ‌های واقعی قطعه نگاشت و مسیریابی شده است. نماهای سطح بالا در زیر آمده‌اند؛ نماهای درونی هر یک از دو پیمانه نیز ثبت شده‌اند، امّا چون بسیار پهن‌اند و در عرض صفحه خوانا نمی‌شوند، تنها در بستهٔ تحویلی آورده شده‌اند.

	\shot{screenshots/post_mapping_top.png}{نمای \lr{Technology Map Viewer (Post-Mapping)} در سطح بالا}
	\shot{screenshots/post_fitting_top.png}{نمای \lr{Technology Map Viewer (Post-Fitting)} در سطح بالا}

	\subsection*{د) تحلیل زمانی}
	برای این پروژه پروندهٔ قیود طراحی\LTRfootnote{Synopsys Design Constraints} تعریف نشده است؛ در چنین حالتی تحلیل‌گر زمانی\LTRfootnote{TimeQuest Timing Analyzer} به‌طور خودکار ساعتی با دورهٔ یک نانوثانیه، معادل یک گیگاهرتز، استنتاج می‌کند. طبیعی است که طرح در برابر چنین قید تحمیلی و غیرواقع‌بینانه‌ای حاشیهٔ زمانی\LTRfootnote{Slack} منفی نشان دهد و هشدار \lr{Timing requirements not met} تولید شود.

	\shot{screenshots/compile_messages.png}{پیام‌های پایانی همگردانی و هشدارهای تحلیل زمانی}

	اعداد گزارش‌شده چنین‌اند: بدترین حاشیهٔ برپایی در مدل تند \(-1.163\) و در مدل کند \(-2.892\) نانوثانیه، و حاشیهٔ نگهداری در هر دو مدل مثبت. همگردانی کامل با صفر خطا به پایان رسید.

	معیار درست ارزیابی در نبود قید، بیشینه‌بسامد\LTRfootnote{Maximum Frequency (Fmax)} گزارش‌شده است. تحلیل‌گر زمانی برای بدترین شرایط کاری مقدار \(256.94\,\mathrm{MHz}\) را گزارش می‌کند و بیشینه‌بسامد مقیّد\LTRfootnote{Restricted Fmax} برابر \(250\,\mathrm{MHz}\) است که از محدودیّت نرخ تغییر پایه‌های ورودی و خروجی\LTRfootnote{Maximum I/O Toggle Rate} قطعه می‌آید و نه از منطق درونی طرح. ساعت ۱۰۰ مگاهرتزی که در میز آزمون به کار رفته، بسیار پایین‌تر از این حدّ است.

	با توجّه به برگزاری مجازی آزمایشگاه، ارزیابی این آزمایش بر پایهٔ شبیه‌سازی، سنتز و تحلیل زمانی نرم‌افزاری انجام شده و در پرونده‌های پیکربندی پروژه تخصیص پایه‌ای برای برد ثبت نشده است؛ ازاین‌رو ادّعایی دربارهٔ آزمون فیزیکی روی برد مطرح نمی‌شود.

	\section*{۷.۱۰ نتایج}
	\begin{itemize}
		\item قاب ده‌بیتی با ترتیب خواسته‌شدهٔ دستورکار، یعنی شروع، توازن، هفت بیت داده و خاتمه، در فرستنده پیاده و در گیرنده بازخوانی شد.
		\item فرستنده و گیرنده هر یک به‌صورت یک ماشین حالت جداگانه نوشته شدند و پیمانهٔ بالا با اتّصال حلقهٔ بازگشتی، مبادلهٔ داده میان آن دو را نشان داد.
		\item طول هر بیت با یک پارامتر بیان شد تا طرح به نرخ باود مشخّصی وابسته نماند.
		\item افزون بر ذخیرهٔ داده، گیرنده سلامت قاب را نیز گزارش می‌کند و هر دو گونهٔ خطای توازن و خطای بیت خاتمه را تشخیص داد.
		\item هر پنج مرحلهٔ میز آزمون خودآزما، شامل دو قاب پیاپی و درخواست هم‌زمان با اشغال، بدون خطا گذشتند.
		\item سنتز با ۲۱۵ عنصر منطقی، ۱۰۶ ثبات و ۲۱ پایه انجام شد و بررسی نشان داد شصت و چهار ثبات از این شمار، تنها به دلیل تعریف شمارنده‌ها از نوع \lr{integer} مصرف شده‌اند.
		\item بیشینه‌بسامد گزارش‌شده \(256.94\,\mathrm{MHz}\) با حدّ مقیّد \(250\,\mathrm{MHz}\) است.
	\end{itemize}

	\section*{نتیجه‌گیری}
	در این آزمایش یک واحد \lr{UART} کامل با فرستنده و گیرندهٔ جداگانه ساخته و درستی مبادلهٔ داده میان آن دو نشان داده شد. آنچه این آزمایش را از یک تمرین ماشین حالت ساده متمایز می‌کند، ناهم‌گام بودن پیوند دو سوی خط است: هیچ ساعت مشترکی میان فرستنده و گیرنده رد و بدل نمی‌شود و تنها قرارداد مشترک، ساختار قاب و طول هر بیت است. به همین دلیل، نمونه‌برداری در میانهٔ هر بیت و نه در لبه‌های آن انجام شد. افزودن وارسی توازن و بیت خاتمه نیز نشان داد که یک قاب کامل‌شده لزوماً یک قاب سالم نیست و این دو مفهوم باید در دو پرچم جداگانه گزارش شوند. سرانجام، بررسی گزارش سنتز یادآور شد که درستی کارکردی و کارآمدی سخت‌افزاری دو چیز متفاوت‌اند؛ چنان‌که در همین طرحِ کاملاً درست، بیش از نیمی از فلیپ‌فلاپ‌ها تنها به سبب یک انتخاب نوع داده مصرف شده‌اند.
\end{document}
